\documentclass[10pt,a4paper]{amsart}
\usepackage{amsmath,amssymb,amsthm,graphicx}
% For commutative diagrams we use Xy-pic
\usepackage[all]{xy}

% We use 2cell for 2-commutative diagrams.
\xyoption{2cell}
% Theorem environments.
%
\theoremstyle{plain}
\newtheorem{theorem}[subsection]{Theorem}
\newtheorem{proposition}[subsection]{Proposition}
\newtheorem{lemma}[subsection]{Lemma}

\theoremstyle{definition}
\newtheorem{definition}[subsection]{Definition}
\newtheorem{example}[subsection]{Example}
\newtheorem{exercise}[subsection]{Exercise}
\newtheorem{situation}[subsection]{Situation}

\theoremstyle{remark}
\newtheorem{remark}[subsection]{Remark}
\newtheorem{remarks}[subsection]{Remarks}

\numberwithin{equation}{subsection}

% Macros
%
\def\lim{\mathop{\mathrm{lim}}\nolimits}
\def\colim{\mathop{\mathrm{colim}}\nolimits}
\def\Spec{\mathop{\mathrm{Spec}}}
\def\Hom{\mathop{\mathrm{Hom}}\nolimits}
\def\Ext{\mathop{\mathrm{Ext}}\nolimits}
\def\SheafHom{\mathop{\mathcal{H}\!\mathit{om}}\nolimits}
\def\SheafExt{\mathop{\mathcal{E}\!\mathit{xt}}\nolimits}
\def\Sch{\mathit{Sch}}
\def\Mor{\mathop{\mathrm{Mor}}\nolimits}
\def\Ob{\mathop{\mathrm{Ob}}\nolimits}
\def\Sh{\mathop{\mathit{Sh}}\nolimits}
\def\NL{\mathop{N\!L}\nolimits}
\def\CH{\mathop{\mathrm{CH}}\nolimits}
\def\proetale{{pro\text{-}\acute{e}tale}}
\def\etale{{\acute{e}tale}}
\def\QCoh{\mathit{QCoh}}
\def\Ker{\mathop{\mathrm{Ker}}}
\def\Im{\mathop{\mathrm{Im}}}
\def\Coker{\mathop{\mathrm{Coker}}}
\def\Coim{\mathop{\mathrm{Coim}}}

% Boxtimes
%
\DeclareMathSymbol{\boxtimes}{\mathbin}{AMSa}{"02}

%
% Macros for moduli stacks/spaces
%
\def\QCohstack{\mathcal{QC}\!\mathit{oh}}
\def\Cohstack{\mathcal{C}\!\mathit{oh}}
\def\Spacesstack{\mathcal{S}\!\mathit{paces}}
\def\Quotfunctor{\mathrm{Quot}}
\def\Hilbfunctor{\mathrm{Hilb}}
\def\Curvesstack{\mathcal{C}\!\mathit{urves}}
\def\Polarizedstack{\mathcal{P}\!\mathit{olarized}}
\def\Complexesstack{\mathcal{C}\!\mathit{omplexes}}
% \Pic is the operator that assigns to X its picard group, usage \Pic(X)
% \Picardstack_{X/B} denotes the Picard stack of X over B
% \Picardfunctor_{X/B} denotes the Picard functor of X over B
\def\Pic{\mathop{\mathrm{Pic}}\nolimits}
\def\Picardstack{\mathcal{P}\!\mathit{ic}}
\def\Picardfunctor{\mathrm{Pic}}
\def\Deformationcategory{\mathcal{D}\!\mathit{ef}}
\usepackage{cleveref}
\crefname{originallemma}{original lemma}{original lemmas}
\Crefname{originallemma}{Original lemma}{Original lemmas}
\makeatletter
\newlabel{definition-gysin-homomorphism}{{02T8}{source}{}{sourcepointer-1}{}}
\newlabel{definition-gysin-homomorphism@cref}{{[originallemma][0][]02T8}{[1][0][]source}{}{}{}}
\makeatother
\crefname{originallemma}{original lemma}{original lemmas}
\Crefname{originallemma}{Original lemma}{Original lemmas}
\makeatletter
\newlabel{lemma-key-formula}{{0AYC}{source}{}{sourcepointer-2}{}}
\newlabel{lemma-key-formula@cref}{{[originallemma][0][]0AYC}{[1][0][]source}{}{}{}}
\makeatother
\crefname{originallemma}{original lemma}{original lemmas}
\Crefname{originallemma}{Original lemma}{Original lemmas}
\makeatletter
\newlabel{remark-infinite-sums-rational-equivalences}{{02RZ}{source}{}{sourcepointer-3}{}}
\newlabel{remark-infinite-sums-rational-equivalences@cref}{{[originallemma][0][]02RZ}{[1][0][]source}{}{}{}}
\makeatother
\crefname{originallemma}{original lemma}{original lemmas}
\Crefname{originallemma}{Original lemma}{Original lemmas}
\makeatletter
\newlabel{section-key}{{0AYB}{source}{}{sourcepointer-4}{}}
\newlabel{section-key@cref}{{[originallemma][0][]0AYB}{[1][0][]source}{}{}{}}
\makeatother
\crefname{originallemma}{original lemma}{original lemmas}
\Crefname{originallemma}{Original lemma}{Original lemmas}
\makeatletter
\newlabel{situation-setup}{{02QL}{source}{}{sourcepointer-5}{}}
\newlabel{situation-setup@cref}{{[originallemma][0][]02QL}{[1][0][]source}{}{}{}}
\makeatother
\title{The Gysin maps attached to two Cartier divisors commute}
\author{The Stacks Project Authors}
\date{stacks-project:chow.tex:local-snapshot-20261012; GNU FDL 1.2 or later}
\begin{document}
\maketitle
\noindent\textit{Editorial note.} All delivered mathematical text is the author's own, copied from the pinned Stacks Project chapter file chow.tex; no statement or proof step was rewritten, reordered or removed.\par\medskip
\noindent\textit{Editorial note.} Preamble provenance: the pinned chapter file opens with \textbackslash{}input\{preamble\} and carries no preamble text of its own, so the project's single Stacks-wide preamble staged at sources/stacks-chow/preamble.tex is excerpted verbatim as two spans, lines 14-18 (Xy-pic setup) and lines 154-217 (the native theorem declarations and mathematical macros). Nothing was added, removed or redefined; the wrapper class and the amsmath/amssymb/amsthm/graphicx base packages are supplied by the assembler.\par\medskip
\noindent\textit{Editorial note.} Same-chapter pointers are located at their \textbackslash{}label lines inside the pinned chapter file; cross-chapter pointers are located in the project's own tag list staged in the snapshot (sources/stacks-chow/tags.txt), at the verbatim lines carrying their canonical tags.\par\medskip
\noindent\textit{Allowed standard prerequisites (editorial).}\begin{itemize}
\item The Gysin homomorphism attached to a triple (invertible sheaf, section, zero scheme). Same chapter chow.tex, label definition-gysin-homomorphism, invoked in the statement at proof line 5500. Fixes what the four arrows in the square are, so that the commuting assertion has a meaning; the comparison itself is carried out in the counted span.
\item The discussion of the key formula: local generators at generic points of divisors. Same chapter chow.tex, label section-key, invoked at proof lines 5534 and 5537. Supplies the normalisation of the meromorphic sections sigma and sigma' at each generic point and the definitions of the multiplicities n\_i and n'\_i, which is the setup of the computation.
\item The key formula and infinite sums of rational equivalences. Same chapter chow.tex, labels lemma-key-formula (proof line 5556) and remark-infinite-sums-rational-equivalences (proof line 5568). Gives the cycle identity turning the difference of the two representatives into a sum of divisors of tame symbols, and lets the resulting rational equivalence pass from the classes [W] to an arbitrary cycle.
\end{itemize}\medskip
\ifdefined\hypertarget\hypertarget{sourcepointer-1}{}\fi
\noindent\textit{Source reference.} Original lemma 02T8: sources/stacks-chow/pinned-original-source.tex, source lines 5069--5069. Navigation pointer for \textbackslash{}ref\{definition-gysin-homomorphism\} used inside the counted span (line 5500); the argument that consumes it is retained in the span.\par\medskip
\ifdefined\hypertarget\hypertarget{sourcepointer-2}{}\fi
\noindent\textit{Source reference.} Original lemma 0AYC: sources/stacks-chow/pinned-original-source.tex, source lines 4746--4746. Navigation pointer for \textbackslash{}ref\{lemma-key-formula\} used inside the counted span (line 5556); the argument that consumes it is retained in the span.\par\medskip
\ifdefined\hypertarget\hypertarget{sourcepointer-3}{}\fi
\noindent\textit{Source reference.} Original lemma 02RZ: sources/stacks-chow/pinned-original-source.tex, source lines 3057--3057. Navigation pointer for \textbackslash{}ref\{remark-infinite-sums-rational-equivalences\} used inside the counted span (line 5568); the argument that consumes it is retained in the span.\par\medskip
\ifdefined\hypertarget\hypertarget{sourcepointer-4}{}\fi
\noindent\textit{Source reference.} Original lemma 0AYB: sources/stacks-chow/pinned-original-source.tex, source lines 4686--4686. Navigation pointer for \textbackslash{}ref\{section-key\} used inside the counted span (line 5534); the argument that consumes it is retained in the span.\par\medskip
\ifdefined\hypertarget\hypertarget{sourcepointer-5}{}\fi
\noindent\textit{Source reference.} Original lemma 02QL: sources/stacks-chow/pinned-original-source.tex, source lines 1372--1372. Navigation pointer for \textbackslash{}ref\{situation-setup\} used inside the counted span (line 5497); the argument that consumes it is retained in the span.\par\medskip
\begin{lemma}
\label{lemma-gysin-commutes-gysin}
Let $(S, \delta)$ be as in Situation \ref{situation-setup}. Let $X$ be locally
of finite type over $S$. Let $(\mathcal{L}, s, i : D \to X)$ and
$(\mathcal{L}', s', i' : D' \to X)$ be two triples as in
Definition \ref{definition-gysin-homomorphism}. Then the diagram
$$
\xymatrix{
\CH_k(X) \ar[r]_{i^*} \ar[d]_{(i')^*} & \CH_{k - 1}(D) \ar[d]^{j^*} \\
\CH_{k - 1}(D') \ar[r]^{(j')^*} & \CH_{k - 2}(D \cap D')
}
$$
commutes where each of the maps is a Gysin map.
\end{lemma}

\begin{proof}
Denote $j : D \cap D' \to D$ and $j' : D \cap D' \to D'$ the closed
immersions corresponding to $(\mathcal{L}|_{D'}, s|_{D'}$ and
$(\mathcal{L}'_D, s|_D)$. We have to show that
$(j')^*i^*\alpha = j^* (i')^*\alpha$ for all $\alpha \in \CH_k(X)$.
Let $W \subset X$ be an integral closed subscheme of dimension $k$.
Let us prove the equality in case $\alpha = [W]$. We will deduce
it from the key formula.

\medskip\noindent
We let $\sigma$ be a nonzero meromorphic section of $\mathcal{L}|_W$
which we require to be equal to $s|_W$ if $W \not \subset D$.
We let $\sigma'$ be a nonzero meromorphic section of $\mathcal{L}'|_W$
which we require to be equal to $s'|_W$ if $W \not \subset D'$.
Write
$$
\text{div}_{\mathcal{L}|_W}(\sigma) =
\sum \text{ord}_{Z_i, \mathcal{L}|_W}(\sigma)[Z_i] = \sum n_i[Z_i]
$$
and similarly
$$
\text{div}_{\mathcal{L}'|_W}(\sigma') =
\sum \text{ord}_{Z_i, \mathcal{L}'|_W}(\sigma')[Z_i] = \sum n'_i[Z_i]
$$
as in the discussion in Section \ref{section-key}.
Then we see that $Z_i \subset D$ if $n_i \not = 0$ and
$Z'_i \subset D'$ if $n'_i \not = 0$. For each $i$, let $\xi_i \in Z_i$
be the generic point. As in Section \ref{section-key} we choose
for each $i$ an element
$\sigma_i \in \mathcal{L}_{\xi_i}$, resp.\ $\sigma'_i \in \mathcal{L}'_{\xi_i}$
which generates over $B_i = \mathcal{O}_{W, \xi_i}$
and which is equal to the image of
$s$, resp.\ $s'$ if $Z_i \not \subset D$, resp.\ $Z_i \not \subset D'$.
Write $\sigma = f_i \sigma_i$ and $\sigma' = f'_i\sigma'_i$ so that
$n_i = \text{ord}_{B_i}(f_i)$ and
$n'_i = \text{ord}_{B_i}(f'_i)$.
From our definitions it follows that
$$
(j')^*i^*[W] =
\sum \text{ord}_{B_i}(f_i) \text{div}_{\mathcal{L}'|_{Z_i}}(\sigma'_i|_{Z_i})
$$
as cycles and
$$
j^*(i')^*[W] =
\sum \text{ord}_{B_i}(f'_i) \text{div}_{\mathcal{L}|_{Z_i}}(\sigma_i|_{Z_i})
$$
The key formula (Lemma \ref{lemma-key-formula}) now gives the equality
$$
\sum \left(
\text{ord}_{B_i}(f_i) \text{div}_{\mathcal{L}'|_{Z_i}}(\sigma'_i|_{Z_i}) -
\text{ord}_{B_i}(f'_i) \text{div}_{\mathcal{L}|_{Z_i}}(\sigma_i|_{Z_i})
\right) =
\sum \text{div}_{Z_i}(\partial_{B_i}(f_i, f'_i))
$$
of cycles. Note that $\text{div}_{Z_i}(\partial_{B_i}(f_i, f'_i)) = 0$ if
$Z_i \not \subset D \cap D'$ because in this case either $f_i = 1$
or $f'_i = 1$. Thus we get a rational equivalence between our specific
cycles representing $(j')^*i^*[W]$ and $j^*(i')^*[W]$ on $D \cap D' \cap W$.
By Remark \ref{remark-infinite-sums-rational-equivalences}
the result follows for general $\alpha$.
\end{proof}

\end{document}
